\documentclass[10pt]{article}
\usepackage[a4paper,margin=0.7in]{geometry}
\usepackage{graphicx,booktabs,amsmath,amssymb,microtype}
\usepackage[hidelinks]{hyperref}
\usepackage{titlesec}
\titlespacing*{\section}{0pt}{6pt}{3pt}
\setlength{\parskip}{3pt}
\setlength{\parindent}{0pt}

\begin{document}
\begin{center}
{\Large \textbf{Comparing Policy Gradient Methods on MuJoCo}}\\[2pt]
CS778: Foundations of Modern AI \quad$\cdot$\quad Roll no.\ 230224
\end{center}

\section{Setup}
VPG, TRPO and PPO (clip) were implemented from scratch in PyTorch and trained on \texttt{HalfCheetah-v4} and \texttt{Walker2d-v4} for 1M environment steps each, with 2 seeds per algorithm--environment pair (12 runs).
To make the comparison fair, \textbf{everything except the policy update is shared}:
\begin{itemize}\setlength\itemsep{0pt}
\item a Gaussian actor (64-64 tanh MLP for the mean, state-independent $\log\sigma$) and a 64-64 tanh critic;
\item running state normalization;
\item GAE advantages ($\gamma=0.99$, $\lambda=0.95$) that bootstrap $V(S_{t+1})$ on time-limit truncation but not on true termination;
\item batches of 2048 steps, and the same critic fit (80 Adam steps on the $\lambda$-return).
\end{itemize}
The three update rules are:
\begin{itemize}\setlength\itemsep{0pt}
\item \textbf{VPG:} one Adam step on $-\frac1N\sum_t \log\pi_\theta(A_t|S_t)\hat A_t$ (lr $3\!\times\!10^{-4}$).
\item \textbf{TRPO:} a natural-gradient step $\sqrt{2\delta/g^\top v}\,v$ with $v\approx F^{-1}g$ from 10 conjugate-gradient iterations on Fisher-vector products, followed by a backtracking line search that requires $\bar D_{KL}\le\delta=0.01$ and surrogate improvement.
\item \textbf{PPO:} the clipped surrogate ($\varepsilon=0.2$), optimized for up to 10 epochs of 64-sample minibatches, with approximate-KL early stopping.
\end{itemize}
Returns below are raw environment returns of the stochastic training policy.

\begin{figure}[h]
\centering
\includegraphics[width=\linewidth]{return_and_length.png}
\vspace{-14pt}
\caption{Episodic return (top) and episode length (bottom) vs.\ environment steps. Each line is the mean of 2 seeds (10-episode moving average) and the band is the min--max over seeds. HalfCheetah cannot terminate, so its length is always 1000.}
\label{fig:main}
\end{figure}

\begin{table}[h]
\centering\small
\begin{tabular}{lrrrr|rrrr}
\toprule
& \multicolumn{4}{c|}{HalfCheetah-v4: return at step} & \multicolumn{4}{c}{Walker2d-v4: return at step} \\
& 100k & 250k & 500k & 1M & 100k & 250k & 500k & 1M \\
\midrule
VPG  & $-128$ & 64   & 326  & 723  & 156 & 231  & 408  & 1202 \\
TRPO & 367    & \textbf{2118} & \textbf{2954} & \textbf{3615} & 338 & 936  & \textbf{2169} & 2983 \\
PPO  & \textbf{549} & 1471 & 1997 & 2761 & 334 & \textbf{1004} & 1829 & \textbf{3628} \\
\bottomrule
\end{tabular}
\caption{Mean over 2 seeds of the smoothed return. Per-seed final returns: HalfCheetah VPG 593/853, TRPO 1614/5604, PPO 2483/3046; Walker2d VPG 1351/932, TRPO 3145/2972, PPO 3458/3910. Walker2d mean length of the last 10 episodes: VPG $\approx$430, TRPO $\approx$760, PPO $\approx$910.}
\label{tab:res}
\end{table}

\section{Observations and explanations}

\textbf{1. Trust-region methods are far more sample-efficient than VPG.}
On both tasks, TRPO and PPO reach in 250k steps what VPG has not reached after 1M (Table~\ref{tab:res}).
VPG uses each batch of 2048 samples for a \emph{single} gradient step.
The step must stay small, because the policy-gradient estimate is only valid near $\theta_k$ and nothing prevents a large step from collapsing the policy.
So each sample buys very little improvement.
TRPO instead takes the \emph{largest} step the KL constraint allows.
Because the step is measured in policy space (through the Fisher matrix), not parameter space, its size adapts automatically to the local geometry.
PPO reuses the same batch for up to 320 minibatch updates, and clipping removes the incentive to push the probability ratio outside $[0.8,1.2]$, so reuse stays safe.
Both methods therefore extract many times more improvement per sample.
This is the practical meaning of treating policy gradient as local optimization of the surrogate $L_{\theta_k}(\theta)$ inside a trust region.

\textbf{2. TRPO vs.\ PPO: comparable, with the ranking decided by seed variance.}
On Walker2d the two track each other for most of training, and PPO is ahead at 1M (3628 vs.\ 2983).
On HalfCheetah, TRPO has the higher mean, but its seeds differ by $\sim$4000.
Watching the trained policies showed that \textbf{TRPO seed 0 flipped the cheetah onto its back} and learned to run upside down: it covers 97\,m per episode against 324\,m for seed 1.
This is a well-known local optimum of HalfCheetah.
A trust region limits \emph{how fast} the policy changes, not \emph{which} optimum it reaches.
Once the policy commits to the flipped gait and its $\log\sigma$ shrinks, no small, KL-bounded step leads back to upright running.
With 2 seeds, TRPO and PPO cannot be separated on HalfCheetah; a flip in one PPO seed would have reversed the ranking.
PPO's seeds happened to stay closer (2483/3046).

\textbf{3. On Walker2d, learning to balance dominates.}
Walker2d gives $+1$ per step while healthy and terminates on a fall, so return and episode length rise together (Fig.~\ref{fig:main}, bottom right).
At 1M steps, VPG walkers fall after $\sim$430 steps.
TRPO and PPO survive $\sim$760 and $\sim$910 steps, and a third to a half of their final episodes reach the 1000-step limit.
Walker2d curves are also much noisier than HalfCheetah's, because a single stumble ends an episode early and cuts its return.
VPG's late jump in return (seed 0 near 550k, seed 1 near 800k steps) coincides with its episodes getting longer, i.e.\ it only then learned to stay upright.

\textbf{4. Stability and cost.}
TRPO's HalfCheetah curves rise almost monotonically, because the line search rejects any step that fails to improve the surrogate inside the KL ball.
PPO's constraint is soft (clipping plus early stopping on an approximate KL), and its curves are slightly noisier.
PPO is much simpler: it needs only first-order gradients, with no Fisher-vector products, conjugate gradient or line search, and it matched TRPO overall.
This matches the lecture's summary: PPO recovers most of TRPO's benefit with first-order machinery.

\textbf{Limitations.}
Two seeds per configuration is enough to show the VPG vs.\ trust-region gap, but not to rank TRPO against PPO.
Hyperparameters were the standard defaults and were not tuned per algorithm.
In particular, VPG with a larger learning rate or a larger batch would likely learn faster, at the cost of more instability.
\end{document}
